\documentclass[12pt]{article}
\usepackage[utf8]{inputenc}
\usepackage[T1]{fontenc}
\usepackage[margin=1in]{geometry}
\usepackage{setspace}
\usepackage{xcolor}
\usepackage[colorlinks=true, urlcolor=blue!60!black, linkcolor=blue!60!black]{hyperref}
\setlength{\parindent}{2em}
\setlength{\parskip}{1.25em}
\renewcommand{\baselinestretch}{1.0}
\begin{document}
\begin{center}
{\LARGE \textit{Zivotofsky v.\ Kerry}, 576 U.S.\ 1 (2015)\par}
\vspace{0.8em}
{\large Case Brief --- Four-Page Edition\par}
\end{center}
\vspace{1em}

\section{Memory Jogger}

The President alone holds the power to recognize foreign sovereigns and their territorial bounds, so Congress could not require the Secretary of State to list ``Israel'' on the passports of Jerusalem-born citizens on request: a presidential victory in \textit{Youngstown}'s lowest-ebb category (\href{https://supreme.justia.com/cases/federal/us/576/1/}{Justia}).

\section{Detailed Case Facts}

Since President Truman recognized Israel in 1948, every President has declined to recognize any nation's sovereignty over Jerusalem, leaving its status to negotiation. The State Department's Foreign Affairs Manual accordingly directed officials to record ``Jerusalem,'' not ``Israel,'' as the birthplace of citizens born there. In 2002, \S~214(d) of the Foreign Relations Authorization Act directed the Secretary, on request, to ``record the place of birth as Israel.'' President George W.\ Bush's signing statement said that a mandatory reading would interfere with the President's authority to determine recognition. Menachem Zivotofsky was born in Jerusalem in 2002 to U.S.\ citizen parents; his mother asked that his passport and consular report list ``Jerusalem, Israel,'' and officials listed only ``Jerusalem.''

\section{Procedural History}

The district court twice dismissed, the second time as a political question, and the D.C.\ Circuit affirmed. In \textit{Zivotofsky I} (2012), the Supreme Court held the case justiciable because it asked only whether a statute is constitutional. On remand, the D.C.\ Circuit held \S~214(d) an unconstitutional infringement on the recognition power. The Supreme Court affirmed on June 8, 2015.

\section{Judicial Votes}

Affirmed: five Justices on the recognition rationale, six on passports.
\begin{itemize}
\item \textbf{Majority:} Justice Kennedy, joined by Justices Ginsburg, Breyer, Sotomayor, and Kagan.
\item \textbf{Concurrences:} Justice Breyer; Justice Thomas (concurring in the judgment in part and dissenting in part).
\item \textbf{Dissents:} Chief Justice Roberts, joined by Alito; Justice Scalia, joined by Roberts and Alito.
\end{itemize}

\section{Holding}

The President has the exclusive power to grant formal recognition to a foreign sovereign and its territorial bounds, and \S~214(d) is unconstitutional because Congress may not command the Executive to issue an official statement in a passport contradicting the President's recognition determination. Because Zivotofsky waived any separate argument about consular reports, the Court addressed only passports.

\section{Analysis of Opinions}

\textbf{Majority (Kennedy).} Under \textit{Youngstown} Category Three, the President prevails only if his power is ``exclusive'' and ``conclusive.'' The recognition power follows from the Reception Clause (the President ``shall receive Ambassadors''), since founding-era jurists treated receiving an ambassador as acknowledging sovereignty, and from the President's role in treaties and appointments. The Nation must ``speak \ldots{} with one voice,'' and only the Executive has ``[d]ecision, activity, secrecy, and dispatch.'' Precedent (\textit{Belmont}, \textit{Pink}, \textit{Sabbatino}) supports exclusivity, but the Court rejected \textit{Curtiss-Wright}'s ``sole organ'' dicta. History, from Citizen Genet to China in 1979, was ``not all on one side,'' but Congress had acquiesced. A passport is a ``political document,'' and \S~214 was designed to force the Secretary to contradict presidential policy.

\textbf{Breyer (concurring).} He still viewed the case as a political question but joined the Court.

\textbf{Thomas (in part).} Listing Israel does not implicate recognition, but the Vesting Clause's residual foreign-affairs power covers passports; consular birth reports fall under the Naturalization Clause, so \S~214(d) is valid as to them.

\textbf{Roberts (dissenting).} The Reception Clause is a duty, not an exclusive grant, and a passport accommodation is not recognition; deferring to foreign misperceptions gives an ``international heckler's veto.''

\textbf{Scalia (dissenting).} \S~214(d) rests on the naturalization power; moving from legal acts to statements is ``a leap worthy of the Mad Hatter,'' and Thomas's theory yields ``a presidency more reminiscent of George III than George Washington.''

\textbf{Later.} \textit{Trump v.\ United States} (2024) cited \textit{Zivotofsky} as an example of ``conclusive and preclusive'' presidential authority.

\section{Examples: Future Applications}

\begin{enumerate}
\item \textbf{Same-Side: Compelled Maps.} A statute requiring official maps and reports to foreign governments that contradict recognition of a disputed province is invalid.
\item \textbf{Same-Side: Receiving a Breakaway Envoy.} Ordering the Secretary to accredit an unrecognized government's ambassador compels the paradigm act of recognition.
\item \textbf{Opposite-Side: Defunding an Embassy.} Congress controls appropriations, though a condition tied to withdrawing recognition invites an unconstitutional-conditions objection.
\item \textbf{Opposite-Side: Sanctions.} Import and arms bans on a recognized government regulate foreign commerce and require no contrary statement.
\item \textbf{Fence-Sitter: Customs Labels.} Requiring customs documents to label goods from an unrecognized ``Independent Republic'' pits control of official statements against the commerce power.
\end{enumerate}

\section{Critique}

\textbf{Formalism splits.} The majority mixes textual inference with ``one voice'' functionalism, while Scalia's enumerated-powers formalism and Thomas's Vesting Clause residualism show originalism fracturing in foreign affairs. \textbf{Statement versus act.} The majority conceded that \S~214(d) ``would not itself constitute a formal act of recognition,'' leaving unclear how far control of ``official executive statements'' reaches. \textbf{History.} Roberts notes that acquiescence normally supports Category Two, not Category Three.

\section{Key Quotations}

\begin{quote}
``This case is confined solely to the exclusive power of the President to control recognition determinations, including formal statements by the Executive Branch acknowledging the legitimacy of a state or government and its territorial bounds. Congress cannot command the President to contradict an earlier recognition determination in the issuance of passports.''
\end{quote}
\noindent --- Justice Kennedy, for the Court, 576 U.S.\ at 32.

\begin{quote}
``The Executive is not free from the ordinary controls and checks of Congress merely because foreign affairs are at issue. \dots\ It is not for the President alone to determine the whole content of the Nation's foreign policy.''
\end{quote}
\noindent --- Justice Kennedy, for the Court, 576 U.S.\ at 21.

\begin{quote}
``Today's decision is a first: Never before has this Court accepted a President's direct defiance of an Act of Congress in the field of foreign affairs.''
\end{quote}
\noindent --- Chief Justice Roberts, dissenting, 576 U.S.\ at 61.

\end{document}
